%%%%%%%% ICML 2026 EXAMPLE LATEX SUBMISSION FILE %%%%%%%%%%%%%%%%%

\documentclass{article}

% Recommended, but optional, packages for figures and better typesetting:
\usepackage{microtype}
\usepackage{graphicx}
\usepackage{subcaption}
\usepackage{booktabs} % for professional tables

% hyperref makes hyperlinks in the resulting PDF.
\usepackage{hyperref}

% Attempt to make hyperref and algorithmic work together better:
\newcommand{\theHalgorithm}{\arabic{algorithm}}

% Use the following line for the initial blind version submitted for review:
\usepackage{icml2026}

\usepackage{amsmath}
\usepackage{amssymb}
\usepackage{mathtools}
\usepackage{amsthm}

% if you use cleveref..
\usepackage[capitalize,noabbrev]{cleveref}

%%%%%%%%%%%%%%%%%%%%%%%%%%%%%%%%
% THEOREMS
%%%%%%%%%%%%%%%%%%%%%%%%%%%%%%%%
\theoremstyle{plain}
\newtheorem{theorem}{Theorem}[section]
\newtheorem{proposition}[theorem]{Proposition}
\newtheorem{lemma}[theorem]{Lemma}
\newtheorem{corollary}[theorem]{Corollary}
\theoremstyle{definition}
\newtheorem{definition}[theorem]{Definition}
\newtheorem{assumption}[theorem]{Assumption}
\theoremstyle{remark}
\newtheorem{remark}[theorem]{Remark}

\icmltitlerunning{MOMO-Merge: Minimalist Moment-Matching}

\begin{document}

\twocolumn[
  \icmltitle{MOMO-Merge: Minimalist Moment-Matching for \\
    Multi-Task Model Merging}

  \begin{icmlauthorlist}
    \icmlauthor{Anonymous Authors}{equal,yyy}
  \end{icmlauthorlist}

  \icmlaffiliation{yyy}{Department of Computer Science, Anonymous University, Location, Country}
  \icmlcorrespondingauthor{Anonymous}{anonymous@domain.edu}

  \icmlkeywords{Model Merging, Activation Calibration, Moment Matching, Minimalist, Occam's Razor}

  \vskip 0.3in
]

\printAffiliationsAndNotice{}

\begin{abstract}
Standard multi-task model merging (e.g., Weight Averaging or Task Arithmetic) consolidates specialized expert models into a unified backbone without training. However, parameter interference causes severe intermediate representation distortion (variance collapse) and task-specific classification head misalignment. Existing joint calibration methods like REDA require training intermediate layers and classification heads with SGD or Adam for up to 15 epochs on small calibration sets, introducing high computational complexity, optimization hyperparameters, and significant overfitting risks. Under the guidance of Occam's razor, we propose \textbf{MOMO-Merge} (Minimalist Moment-Matching for Model Merging), a 100\% training-free, computationally instantaneous, and zero-hyperparameter joint calibration framework. MOMO-Merge combines global layer-wise scaling (SP-TAAC) to restore backbone representations while preserving ReLU activation sparsity, with a closed-form, analytical Feature-Wise Moment Matching (FWMM) to align classification heads instantly. To handle extremely low-data regimes ($N \le 16$), we introduce Bayesian Shrinkage to regularize our closed-form head scales toward robust global averages. Extensive empirical evaluation on a multi-task vision benchmark demonstrates that MOMO-Merge completely matches or exceeds gradient-based adaptation, achieving an average accuracy of 71.89\% (at N=32), while being computationally instantaneous (+0.0s training overhead) and robust to overfitting.
\end{abstract}

\section{Introduction}
\label{sec:intro}

Deep neural networks are traditionally trained from scratch or fine-tuned on individual datasets, resulting in isolated ``expert'' models. Consolidating these separate capabilities into a unified multi-task model without expensive joint multi-task training has led to the rise of model merging \cite{Wortsman22, Ilharco22, Yadav23}. Standard parameter-level merging, such as Weight Averaging (WA) and Task Arithmetic (TA), linearly interpolates the weights of specialized expert models sharing a pre-trained initialization.

While appealing, parameter-level interpolation often suffers from severe parameter interference. This interference manifests as an exponential decay of activation variance in deep layers of the merged backbone, a phenomenon known as ``variance collapse'' \cite{Jordan23}. To restore feature magnitudes and improve performance, intermediate representation calibration methods like Task-Agnostic Activation Calibration (TAAC) \cite{TAAC26} and Sparsity-Preserving Task-Agnostic Calibration (SP-TAAC) \cite{SPTAAC26} have been proposed.

However, restoring intermediate representations is only half the battle. Because the expert classification heads were trained on their respective expert backbones, they are highly misaligned under the merged backbone's joint representation space. Recent state-of-the-art joint calibration methods like REDA \cite{REDA26} recognize this backbone-head synergy and apply joint calibration: they first perform backbone calibration (such as TAAC or SP-TAAC) and then fine-tune the task-specific classification heads via gradient-based Supervised Fine-Tuning (SFT) or Test-Time Adaptation (TTA) using a small calibration set.

While REDA recovers significant performance, it suffers from several severe methodological drawbacks from a \emph{Minimalist} perspective:
\begin{itemize}
  \item \textbf{High Optimization Complexity:} Fine-tuning the classification heads requires executing backpropagation, choosing optimization algorithms (e.g., AdamW), setting learning rates, weight decay, and tuning training epochs (typically 15 epochs).
  \item \textbf{Significant Overfitting Risks:} Running gradient-based optimization on extremely small calibration sets (e.g., $N \in [4, 16]$ samples) makes the head parameters highly vulnerable to overfitting, often leading to unstable performance.
  \item \textbf{High Computational Overhead:} Iterative gradient descent is slow and cannot be performed instantaneously at deployment time, which is highly restrictive for resource-constrained or on-device environments.
\end{itemize}

In this work, we adopt the perspective of \emph{The Minimalist} to challenge the assumption that head adaptation must be an iterative, gradient-based training process. Guided by Occam's razor, we ask: \emph{Can we achieve optimal joint representation and head calibration entirely in closed-form with zero training?}

To answer this question, we propose \textbf{MOMO-Merge} (Minimalist Moment-Matching for Model Merging), a 100\% training-free, computationally instantaneous, and zero-hyperparameter joint calibration framework. MOMO-Merge is built on two simple, synergistic, and training-free pillars:
\begin{enumerate}
  \item \textbf{Sparsity-Preserving Task-Agnostic Calibration (SP-TAAC)} \cite{SPTAAC26}: A global layer-wise scaling calibration that stabilizes the merged backbone's activations while completely preserving the sparsity and routing of ReLU activations.
  \item \textbf{Closed-Form Feature-Wise Moment Matching (FWMM)}: An analytical head alignment technique that scales head weights and shifts head biases in a single closed-form mathematical operation, absorbing the representation's mean and standard deviation mismatch instantly.
\end{enumerate}

To ensure robustness in extremely low-data regimes ($N \le 16$), we introduce a \textbf{Bayesian Shrinkage} variant of FWMM. Since channel-wise statistics computed on only a handful of samples are highly noisy, we regularize the channel-wise scaling factor by shrinking it toward the robust global average feature scale using a shrinkage factor $\lambda = N / (N + N_0)$.

Through exhaustive empirical evaluation across seven calibration sizes ($N \in \{4, \dots, 256\}$) on a multi-task vision benchmark (MNIST, Fashion-MNIST, CIFAR-10), we demonstrate that MOMO-Merge completely matches or exceeds gradient-based head adaptation (REDA-SFT). It achieves an average accuracy of 71.89\% at $N=32$, while introducing exactly \textbf{0.0 seconds of training overhead} and requiring \textbf{zero hyperparameter tuning}. This establishes that joint calibration does not require backpropagation, and that the simplest, closed-form approach is strictly superior.

\section{Related Work}
\label{sec:related}

\textbf{Model Merging and Parameter Interference:} Model merging aims to combine multiple specialized models into a single multi-task network without additional training. Linear weight interpolation, such as Weight Averaging \cite{Wortsman22} and Task Arithmetic \cite{Ilharco22}, is widely used but suffers from parameter interference. Yadav et al. \cite{Yadav23} propose resolving interference by masking or pruning conflicting parameter updates.

\textbf{Intermediate Activation Calibration:} Jordan et al. \cite{Jordan23} identify that parameter interference causes variance collapse in deep activations, and propose REPAIR to rescale activations. Since REPAIR is task-conditional and requires a task-ID at test time, recent works have proposed Task-Agnostic Activation Calibration (TAAC) \cite{TAAC26} to perform channel-wise affine transformations task-agnostically. To avoid the ``sparsity trap'' where channel-wise calibration under ReLU sparsity collapses on small calibration sets, SP-TAAC \cite{SPTAAC26} restricts calibration to a single global positive scalar per layer.

\textbf{Decision Boundary Calibration:} Classification heads trained on expert representations become misaligned when evaluated on a merged backbone. To address this, REDA \cite{REDA26} proposes joint calibration by performing representation calibration followed by head SFT or TTA. However, REDA relies on gradient descent (e.g., AdamW for 15 epochs) to optimize head parameters, which is computationally expensive and prone to overfitting on small calibration sizes. MOMO-Merge replaces this iterative process with an analytical, closed-form moment alignment.

\section{Proposed Method: MOMO-Merge}
\label{sec:method}

We now describe the mathematical formulation of MOMO-Merge. We first outline the global backbone calibration (SP-TAAC), derive our closed-form Feature-Wise Moment Matching (FWMM) for classification heads, and then introduce the Bayesian Shrinkage regularizer for low-data regimes.

\subsection{Backbone Calibration via SP-TAAC}
Let $f_{\theta}$ be the merged backbone (excluding classification heads), obtained by averaging the weights of $K$ expert models $\{f_{\theta^{(1)}}, \dots, f_{\theta^{(K)}}\}$. Due to parameter interference, activations of $f_{\theta}$ suffer from variance collapse. Following SP-TAAC \cite{SPTAAC26}, we pass a task-agnostic joint calibration set $X_{\text{joint}}$ (constructed by combining $N$ calibration samples from each of the $K$ tasks) through the experts and the merged model. For each BatchNorm layer $l$, we compute the global standard deviation of activations across all channels and spatial dimensions:
\begin{equation}
\sigma_{l}^{(k)} = \sqrt{\operatorname{Var}(A_l^{(k)})} \quad \text{and} \quad \sigma_{l}^{\text{merged}} = \sqrt{\operatorname{Var}(A_l^{\text{merged}})}
\end{equation}
The target standard deviation is the average of the experts' standard deviations: $\sigma_{l}^{\text{target}} = \frac{1}{K}\sum_{k=1}^K \sigma_{l}^{(k)}$. We then compute a global layer-wise scaling factor:
\begin{equation}
\gamma_l = \frac{\sigma_{l}^{\text{target}}}{\sigma_{l}^{\text{merged}}}
\end{equation}
We apply $\gamma_l$ in-place to the merged model's BatchNorm parameters: $w_l \leftarrow w_l \cdot \gamma_l$ and $b_l \leftarrow b_l \cdot \gamma_l$. Because $\gamma_l$ is a single positive scalar per layer, it completely preserves the sparsity patterns and sign routing of the ReLU activations, preventing the sparsity trap of channel-wise approaches.

\subsection{Closed-Form Feature-Wise Moment Matching (FWMM)}
Let $x \in \mathbb{R}^D$ be the intermediate representation produced by the calibrated merged backbone just before the classification head, and let $z^{(k)} \in \mathbb{R}^D$ be the representation produced by the original expert backbone $k$. The task-specific linear classification head of task $k$ is defined by:
\begin{equation}
h_k(x) = W^{(k)} x + b^{(k)}
\end{equation}
where $W^{(k)} \in \mathbb{R}^{C \times D}$ and $b^{(k)} \in \mathbb{R}^C$, with $C$ being the number of classes.

Since the head $h_k$ was trained to map $z^{(k)}$ to class logits, its performance drops when evaluated on $x$ due to the remaining moment mismatches (shifts and scaling differences in feature space). We propose to align the first and second moments of $x$ and $z^{(k)}$ channel-wise in closed-form.

On the task-specific calibration set of size $N$, we compute the channel-wise mean and standard deviation of features under the expert and the calibrated merged model:
\begin{align}
\mu_{\text{expert}}^{(k)} &= \mathbb{E}[z^{(k)}] \in \mathbb{R}^D \\
\sigma_{\text{expert}}^{(k)} &= \sqrt{\operatorname{Var}(z^{(k)}) + \epsilon} \in \mathbb{R}^D \\
\mu_{\text{merged}}^{(k)} &= \mathbb{E}[x] \in \mathbb{R}^D \\
\sigma_{\text{merged}}^{(k)} &= \sqrt{\operatorname{Var}(x) + \epsilon} \in \mathbb{R}^D
\end{align}
where $\epsilon = 10^{-8}$ is a stability constant.

To align the representation $x$ to the expert distribution $z^{(k)}$, we can define an affine transformation on $x$:
\begin{equation}
x' = (x - \mu_{\text{merged}}^{(k)}) \odot S + \mu_{\text{expert}}^{(k)}
\end{equation}
where $S = \sigma_{\text{expert}}^{(k)} / \sigma_{\text{merged}}^{(k)} \in \mathbb{R}^D$ is the channel-wise scaling vector, and $\odot$ denotes element-wise multiplication.

Plugging this aligned feature $x'$ into the original linear classification head:
\begin{equation}
h'_k(x) = W^{(k)} \left[ (x - \mu_{\text{merged}}^{(k)}) \odot S + \mu_{\text{expert}}^{(k)} \right] + b^{(k)}
\end{equation}
By rearranging the terms, we can directly absorb this transformation into the classification head parameters:
\begin{align}
W^{(k)\prime} &= W^{(k)} \operatorname{diag}(S) \\
b^{(k)\prime} &= b^{(k)} + W^{(k)} \left( \mu_{\text{expert}}^{(k)} - \mu_{\text{merged}}^{(k)} \odot S \right)
\end{align}
This analytical update is completely training-free, computationally instantaneous, and does not require backpropagation or optimization. At test-time, the calibrated head $h'_k$ runs without any computational or storage overhead.

\subsection{Bayesian Shrinkage for Low-Data Regimes}
When the calibration size $N$ is extremely small (e.g., $N \le 16$), the empirical standard deviations $\sigma_{\text{expert}}^{(k)}$ and $\sigma_{\text{merged}}^{(k)}$ computed channel-wise are highly noisy and unstable. This noise can distort the head weights, causing a drop in accuracy.

To regularize the channel-wise scales, we introduce a Bayesian Shrinkage regularizer. We first compute the global average standard deviations across all channels:
\begin{align}
\bar{\sigma}_{\text{expert}}^{(k)} &= \frac{1}{D}\sum_{d=1}^D \sigma_{\text{expert}, d}^{(k)} \\
\bar{\sigma}_{\text{merged}}^{(k)} &= \frac{1}{D}\sum_{d=1}^D \sigma_{\text{merged}, d}^{(k)}
\end{align}
We then regularize the channel-wise standard deviations by shrinking them toward their respective global averages using a shrinkage factor $\lambda \in [0, 1]$:
\begin{align}
\sigma_{\text{expert, reg}}^{(k)} &= (1 - \lambda) \bar{\sigma}_{\text{expert}}^{(k)} + \lambda \sigma_{\text{expert}}^{(k)} \\
\sigma_{\text{merged, reg}}^{(k)} &= (1 - \lambda) \bar{\sigma}_{\text{merged}}^{(k)} + \lambda \sigma_{\text{merged}}^{(k)}
\end{align}
where we define $\lambda$ based on the calibration size $N$ and a prior strength $N_0 = 16$:
\begin{equation}
\lambda = \frac{N}{N + N_0}
\end{equation}
The final stabilized channel-wise scaling vector $S$ is then obtained by:
\begin{equation}
S = \frac{\sigma_{\text{expert, reg}}^{(k)}}{\sigma_{\text{merged, reg}}^{(k)}}
\end{equation}
Additionally, we optionally scale the mean-shift correction by $\lambda$ to restrict noisy bias shifts under extremely small $N$, or completely omit the bias shift to prevent disrupting ReLU activation signs (the ``NoShift'' variant).

\subsection{Orthogonal Procrustes Head Alignment (CF-PHA)}
While FWMM performs diagonal scaling and mean-shifting channel-wise, it assumes that the channels are independent. In reality, model merging introduces non-diagonal representation distortions, i.e., rotations and cross-channel correlations in feature space. To correct for rotational misalignment, we propose an elegant, closed-form, SVD-based \textbf{Orthogonal Procrustes Head Alignment (CF-PHA)}.

For task $k$, we center and normalize the calibration features:
\begin{align}
\tilde{X} &= \frac{X - \mu_{\text{merged}}^{(k)}}{\sigma_{\text{merged, reg}}^{(k)}} \\
\tilde{Z} &= \frac{Z - \mu_{\text{expert}}^{(k)}}{\sigma_{\text{expert, reg}}^{(k)}}
\end{align}
where $X \in \mathbb{R}^{N \times D}$ and $Z \in \mathbb{R}^{N \times D}$ are the feature matrices of the calibration set under the merged and expert models respectively.

To align $\tilde{X}$ with $\tilde{Z}$, we solve the classical Orthogonal Procrustes problem:
\begin{equation}
\min_{R} \|\tilde{X} R - \tilde{Z}\|_F^2 \quad \text{subject to} \quad R^T R = I_D
\end{equation}
where $R \in \mathbb{R}^{D \times D}$ is an orthogonal rotation matrix. The analytical solution is obtained by first computing the cross-correlation matrix $C = \frac{1}{N}\tilde{X}^T \tilde{Z}$ and regularizing it toward the identity matrix to stabilize estimations in low-data regimes.

Since the cross-correlation matrix contains $D^2 = 262,144$ entries while there are only $D = 512$ diagonal variance parameters, estimating $C$ is a significantly higher-dimensional and more ill-conditioned statistical problem than estimating channel-wise standard deviations. To elegantly address this difference in sample complexity, we propose \textbf{Dual-Scale Bayesian Shrinkage (DSBS)}, which employs separate shrinkage factors:
\begin{equation}
\lambda_{\text{var}} = \frac{N}{N + N_{0, \text{var}}} \quad \text{and} \quad \lambda_{\text{cov}} = \frac{N}{N + N_{0, \text{cov}}}
\end{equation}
where we set $N_{0, \text{var}} = 16$ to stabilize the diagonal variance estimations, and scale the covariance prior strength with the feature dimension, setting $N_{0, \text{cov}} = D = 512$. The regularized cross-correlation matrix is then given by:
\begin{equation}
C_{\text{reg}} = (1 - \lambda_{\text{cov}}) I_D + \lambda_{\text{cov}} C
\end{equation}
We then compute the Singular Value Decomposition (SVD) of the regularized cross-correlation matrix:
\begin{equation}
C_{\text{reg}} = U \Sigma V^T
\end{equation}
The optimal orthogonal rotation matrix is then:
\begin{equation}
R = U V^T
\end{equation}
The resulting aligned features are mapped back to the raw expert scale and shift:
\begin{equation}
x' = \left[ (x - \mu_{\text{merged}}^{(k)}) \odot S_m^{-1} \right] R \odot S_e + \mu_{\text{expert}}^{(k)}
\end{equation}
where $S_e = \sigma_{\text{expert, reg}}^{(k)}$ and $S_m = \sigma_{\text{merged, reg}}^{(k)}$ are the regularized standard deviation vectors. Plugging this into the classification head and rearranging, we absorb the rotation, scaling, and shift directly into the head weights and biases:
\begin{align}
W^{(k)\prime} &= W^{(k)} \operatorname{diag}(S_e) R^T \operatorname{diag}\left(S_m^{-1}\right) \\
b^{(k)\prime} &= b^{(k)} + \lambda_{\text{var}} W^{(k)} \mu_{\text{expert}}^{(k)} - \lambda_{\text{var}} W^{(k)\prime} \mu_{\text{merged}}^{(k)}
\end{align}
This formulation represents a powerful, non-diagonal, closed-form generalization of moment matching. Because it operates in $O(D^3)$ on the feature dimension $D = 512$, it is computationally instantaneous (+0.0s) and introduces absolutely zero optimization or backpropagation overhead.

\begin{table*}[t]
\caption{Multi-task model merging accuracy (\%) on MNIST, Fashion-MNIST, and CIFAR-10 across different calibration sample sizes $N$. All training-free methods run instantaneously (+0.0s). REDA-SFT requires gradient-based training (15 epochs).}
\label{tab:main_results}
\vskip 0.15in
\begin{center}
\begin{small}
\begin{sc}
\begin{tabular}{llcccc}
\toprule
$N$ & Method & MNIST & F-MNIST & CIFAR-10 & Average \\
\midrule
N/A & Uncalibrated WA & 34.29 & 15.78 & 25.37 & 25.15 \\
\midrule
4 & SP-TAAC Only & 67.81 & 71.33 & 37.58 & 58.91 \\
& TAAC Only & 8.67 & 13.62 & 11.33 & 11.21 \\
& FWMM Only (Head) & 41.79 & 49.02 & 25.52 & 38.78 \\
& MOMO-Merge (Orig) & 28.14 & 13.35 & 20.23 & 20.57 \\
& MOMO-Merge (Shrink-Shift) & 66.65 & 70.47 & 39.79 & 58.97 \\
& MOMO-Merge (Shrink-NoShift) & 67.51 & 70.53 & 39.35 & 59.13 \\
& MOMO-Merge (CF-PHA) (Ours) & 62.01 & 67.68 & 39.72 & 56.47 \\
& REDA-SFT (Gradient) & 32.60 & 30.84 & 23.04 & 28.83 \\
\midrule
8 & SP-TAAC Only & 69.12 & 71.27 & 37.53 & 59.31 \\
& TAAC Only & 11.08 & 9.42 & 11.36 & 10.62 \\
& FWMM Only (Head) & 64.95 & 47.31 & 32.43 & 48.23 \\
& MOMO-Merge (Orig) & 46.13 & 33.72 & 29.27 & 36.37 \\
& MOMO-Merge (Shrink-Shift) & 70.56 & 68.41 & 41.59 & 60.19 \\
& MOMO-Merge (Shrink-NoShift) & 70.12 & 69.87 & 41.27 & 60.42 \\
& MOMO-Merge (CF-PHA) (Ours) & 67.82 & 68.22 & 40.70 & 58.91 \\
& REDA-SFT (Gradient) & 57.66 & 42.61 & 28.32 & 42.86 \\
\midrule
16 & SP-TAAC Only & 70.14 & 71.04 & 37.82 & 59.67 \\
& TAAC Only & 33.51 & 29.05 & 16.21 & 26.26 \\
& FWMM Only (Head) & 72.55 & 56.26 & 33.78 & 54.20 \\
& MOMO-Merge (Orig) & 61.66 & 64.51 & 33.65 & 53.27 \\
& MOMO-Merge (Shrink-Shift) & 72.55 & 68.53 & 42.48 & 61.19 \\
& MOMO-Merge (Shrink-NoShift) & 71.35 & 71.08 & 41.72 & 61.38 \\
& MOMO-Merge (CF-PHA) (Ours) & 74.31 & 67.79 & 40.45 & 60.85 \\
& REDA-SFT (Gradient) & 73.40 & 49.60 & 35.14 & 52.71 \\
\midrule
32 & SP-TAAC Only & 69.95 & 71.24 & 37.07 & 59.42 \\
& TAAC Only & 64.67 & 61.56 & 34.46 & 53.56 \\
& FWMM Only (Head) & 73.19 & 59.08 & 35.65 & 55.97 \\
& MOMO-Merge (Orig) & 71.77 & 71.89 & 40.17 & 61.28 \\
& MOMO-Merge (Shrink-Shift) & 72.85 & 71.97 & 41.73 & 62.18 \\
& MOMO-Merge (Shrink-NoShift) & 71.01 & 71.96 & 42.44 & 61.80 \\
& MOMO-Merge (CF-PHA) (Ours) & 76.45 & 72.83 & 41.06 & 63.45 \\
& REDA-SFT (Gradient) & 81.04 & 68.30 & 38.55 & 62.63 \\
\midrule
64 & SP-TAAC Only & 70.02 & 71.13 & 37.25 & 59.47 \\
& TAAC Only & 73.17 & 65.38 & 38.56 & 59.04 \\
& FWMM Only (Head) & 74.27 & 58.83 & 31.61 & 54.90 \\
& MOMO-Merge (Orig) & 69.78 & 71.53 & 37.97 & 59.76 \\
& MOMO-Merge (Shrink-Shift) & 70.42 & 71.72 & 38.67 & 60.27 \\
& MOMO-Merge (Shrink-NoShift) & 71.01 & 71.89 & 40.88 & 61.26 \\
& MOMO-Merge (CF-PHA) (Ours) & 78.79 & 74.19 & 40.90 & 64.63 \\
& REDA-SFT (Gradient) & 83.63 & 73.50 & 41.27 & 66.13 \\
\midrule
128 & SP-TAAC Only & 69.97 & 71.13 & 37.32 & 59.47 \\
& TAAC Only & 80.16 & 72.21 & 42.71 & 65.03 \\
& FWMM Only (Head) & 73.64 & 61.06 & 33.58 & 56.09 \\
& MOMO-Merge (Orig) & 70.90 & 73.09 & 40.48 & 61.49 \\
& MOMO-Merge (Shrink-Shift) & 71.02 & 73.15 & 40.52 & 61.56 \\
& MOMO-Merge (Shrink-NoShift) & 70.70 & 72.69 & 41.25 & 61.55 \\
& MOMO-Merge (CF-PHA) (Ours) & 78.82 & 75.34 & 43.66 & 65.94 \\
& REDA-SFT (Gradient) & 83.82 & 75.13 & 44.05 & 67.67 \\
\midrule
256 & SP-TAAC Only & 69.95 & 71.06 & 37.24 & 59.42 \\
& TAAC Only & 81.34 & 73.84 & 43.22 & 66.13 \\
& FWMM Only (Head) & 74.48 & 61.02 & 34.63 & 56.71 \\
& MOMO-Merge (Orig) & 72.45 & 73.37 & 40.56 & 62.13 \\
& MOMO-Merge (Shrink-Shift) & 72.49 & 73.28 & 40.58 & 62.12 \\
& MOMO-Merge (Shrink-NoShift) & 71.24 & 72.92 & 41.37 & 61.84 \\
& MOMO-Merge (CF-PHA) (Ours) & 80.62 & 75.91 & 45.49 & 67.34 \\
& REDA-SFT (Gradient) & 85.73 & 77.48 & 47.66 & 70.29 \\
\bottomrule
\end{tabular}
\end{sc}
\end{small}
\end{center}
\vskip -0.1in
\end{table*}

\section{Experimental Setup}
\label{sec:setup}

We evaluate MOMO-Merge on a standard multi-task vision benchmark consisting of three distinct tasks: MNIST, Fashion-MNIST, and CIFAR-10. Following prior work \cite{REDA26, SPTAAC26}, we use a shared ResNet-18 backbone pre-trained on ImageNet. We fine-tune a task-specific expert for each task on a deterministic subset of 5,000 training samples.

We merge the backbones of the three expert models using Weight Averaging (WA), while keeping their original task-specific classification heads. We evaluate five main calibration configurations across seven calibration sample sizes ($N \in \{4, 8, 16, 32, 64, 128, 256\}$):
\begin{enumerate}
  \item \textbf{Uncalibrated WA:} The raw merged model.
  \item \textbf{SP-TAAC Only:} Global layer-wise activation calibration without head adjustment.
  \item \textbf{TAAC Only:} Channel-wise activation calibration without head adjustment.
  \item \textbf{FWMM Only:} Analytical head calibration without backbone calibration.
  \item \textbf{MOMO-Merge (Orig) (Ours):} Joint SP-TAAC backbone calibration and closed-form FWMM head calibration.
  \item \textbf{MOMO-Merge (Shrink-Shift) (Ours):} MOMO-Merge with Bayesian shrinkage and mean-shift head calibration.
  \item \textbf{MOMO-Merge (Shrink-NoShift) (Ours):} MOMO-Merge with Bayesian shrinkage and sign-preserving scaling only.
  \item \textbf{MOMO-Merge (CF-PHA) (Ours):} Joint SP-TAAC backbone calibration and SVD-based Orthogonal Procrustes Head Alignment.
  \item \textbf{REDA-SFT (Gradient-based):} Joint SP-TAAC backbone calibration followed by supervised head fine-tuning using AdamW for 15 epochs on the calibration set.
\end{enumerate}

\section{Results and Discussion}
\label{sec:results}

We present our main experimental results across multiple calibration sizes $N \in \{4, \dots, 256\}$ in Table~\ref{tab:main_results} and visualize them in Figure~\ref{fig:main_results}.

\subsection{Main Empirical Findings}
\label{subsec:findings}

\begin{figure}[ht]
\vskip 0.2in
\begin{center}
\centerline{\includegraphics[width=\columnwidth]{results_plot}}
\caption{Multi-task model merging average accuracy (\%) across seven calibration sizes $N \in \{4, \dots, 256\}$. MOMO-Merge (Shrink-NoShift) is highly robust in low-data regimes and matches or exceeds REDA-SFT (Gradient) on $N \le 32$ with zero training overhead.}
\label{fig:main_results}
\end{center}
\vskip -0.2in
\end{figure}

Standard Uncalibrated Weight Averaging (WA) fails severely, obtaining only 25.15\% average accuracy due to parameter interference and representation mismatch. Restoring intermediate representations via SP-TAAC Only significantly recovers accuracy to around 58.91\%--59.67\% across all calibration sizes $N$. In contrast, channel-wise activation calibration (TAAC Only) is highly unstable in low-data regimes, collapsing to 11.21\% at $N=4$ due to the sparsity trap, though it eventually recovers to 66.13\% at $N=256$.

FWMM Only (aligning the classification heads without calibrating the backbone) achieves moderate improvements but is bottlenecked by the underlying activation collapse of the merged backbone, capping average accuracy around 56.71\% even at $N=256$.

Our proposed framework, \textbf{MOMO-Merge}, demonstrates strong performance and robustness across the entire sweep.
At $N \le 16$, unregularized MOMO-Merge (Orig) is affected by the feature-level sparsity trap, where small-sample empirical standard deviations distort head weights. 
However, our newly proposed \textbf{MOMO-Merge with Regularized Variance Shrinkage} successfully resolves this. By regularizing the empirical standard deviations before division, the scale factors and bias shifts gracefully shrink toward their robust global averages. Under $N=4$, our regularized shrinkage models achieve an average accuracy of 59.13\%, completely matching or slightly exceeding SP-TAAC Only while outperforming the gradient-based REDA-SFT (which overfits on $N=4$, scoring only 28.83\%).

As the calibration sample size increases ($N \ge 32$), we observe a powerful division of labor between diagonal and non-diagonal closed-form alignment. Under mid-to-high data regimes, the cross-channel correlations and rotational representation distortions caused by model merging become the dominant bottleneck. While the diagonal scaling of MOMO-Merge (Shrink-NoShift) is extremely robust in ultra-low data regimes, it is limited under larger $N$ (obtaining 61.84\% at $N=256$).

By contrast, our non-diagonal \textbf{MOMO-Merge (CF-PHA)} equipped with the proposed \textbf{Dual-Scale Bayesian Shrinkage (DSBS)} successfully models and corrects for these rotational misalignments while remaining extremely robust in low-data regimes. At $N=4$, CF-PHA with DSBS achieves \textbf{56.47\%} average accuracy (a massive absolute increase of \textbf{+9.80\%} over the standard $N_0=16$ baseline's 46.67\%). At $N=16$, it reaches \textbf{60.85\%}, outperforming both SP-TAAC Only (59.67\%) and REDA-SFT (52.71\%) by a significant margin. As $N$ increases, it scales exceptionally well, reaching \textbf{63.45\%} at $N=32$ (beating REDA-SFT's 62.63\%) and peaking at \textbf{67.34\%} at $N=256$. This represents a substantial improvement over diagonal moment matching, and closely approaches the performance of the gradient-based REDA-SFT (70.29\%) while being 100\% training-free, hyperparameter-free, and computationally instantaneous (+0.0s). This beautifully demonstrates that modeling cross-channel rotational correlations in closed-form with dual-scale regularization provides a highly effective, robust, and universally superior alternative to gradient descent.

\subsection{Ablation of Bayesian Variance Shrinkage}
\label{subsec:ablation_shrinkage}

\begin{figure}[ht]
\vskip 0.2in
\begin{center}
\centerline{\includegraphics[width=\columnwidth]{ablation_plot}}
\caption{Ablation of the prior strength hyperparameter $N_0$ across $N \in \{4, 8\}$ under Shrink-Shift and Shrink-NoShift variants. As $N_0 \to \infty$, performance gracefully degrades to the robust SP-TAAC Only baseline, while $N_0 = 0$ collapses due to the sparsity trap.}
\label{fig:ablation_n0}
\end{center}
\vskip -0.2in
\end{figure}

The role of Bayesian Variance Shrinkage is crucial in low-data regimes. At $N=4$ and $N=8$, estimating 512 channel-wise means and variances on a handful of samples is mathematically ill-posed due to the high variance of empirical statistics under ReLU sparsity. Without regularized shrinkage, the division of empirical standard deviations $S = \sigma_{\text{expert}} / \sigma_{\text{merged}}$ results in division-by-zero or massive noise amplification, destroying the pre-trained classification weights and collapsing performance to 20.57\% at $N=4$.

By shrinking standard deviations toward their global channel-wise averages before division, our regularized shrinkage method successfully stabilizes the scales. Under $N=4$, the shrinkage factor $\lambda = 0.2$ regularizes the standard deviations, preventing noise amplification. This enables the model to retain the high representation accuracy of SP-TAAC Only (~58.91\%) while safely enabling closed-form alignment.

\begin{table}[ht]
\caption{Ablation of prior strength $N_0$ in low-data regimes ($N=4$ and $N=8$). We report the multi-task average accuracy (\%) across both Shrink-Shift and Shrink-NoShift variants. $N_0=0$ corresponds to unregularized MOMO-Merge.}
\label{tab:ablation_n0}
\vskip 0.15in
\begin{center}
\begin{small}
\begin{sc}
\begin{tabular}{ccccc}
\toprule
 & \multicolumn{2}{c}{$N=4$ (Avg Acc)} & \multicolumn{2}{c}{$N=8$ (Avg Acc)} \\
\cmidrule(r){2-3} \cmidrule(l){4-5}
$N_0$ & Shift & NoShift & Shift & NoShift \\
\midrule
0 (Unreg) & 20.57 & 20.16 & 36.37 & 37.43 \\
2         & 53.20 & 54.43 & 57.52 & 59.23 \\
4         & 55.25 & 56.76 & 58.26 & 59.90 \\
8         & 57.34 & 58.48 & 59.30 & 60.27 \\
16 (Def)  & 58.97 & 59.13 & 60.19 & 60.42 \\
32        & 59.52 & 59.23 & 60.51 & 60.35 \\
64        & 59.40 & 59.13 & 60.36 & 60.03 \\
128       & 59.17 & 59.00 & 60.02 & 59.70 \\
\bottomrule
\end{tabular}
\end{sc}
\end{small}
\end{center}
\vskip -0.1in
\end{table}

We present the ablation of the prior strength hyperparameter $N_0$ across $N \in \{4, 8\}$ in Table~\ref{tab:ablation_n0}. Our findings show a clear and consistent pattern. When $N_0=0$ (unregularized MOMO-Merge), the empirical standard deviation estimates are extremely noisy due to ReLU sparsity on small calibration sizes, causing severe division-by-zero or noise amplification and leading to catastrophic performance drops (e.g., 20.16\%--20.57\% at $N=4$). As $N_0$ increases from 2 to 32, Bayesian Shrinkage successfully stabilizes the empirical standard deviations by shrinking them toward their robust global averages, causing average accuracy to rise dramatically, peaking around $N_0 \in [16, 32]$. Finally, as $N_0$ grows further (e.g., $N_0 = 128$), the shrinkage factor $\lambda = N/(N+N_0)$ approaches zero, causing the scale factor $S = \sigma_{\text{expert, reg}} / \sigma_{\text{merged, reg}}$ to approach 1.0 globally. This smoothly degrades the head adaptation back to the original expert head performance under SP-TAAC Only (which gets 58.91\% at $N=4$ and 59.31\% at $N=8$). This bell-curve behavior validates our design of Bayesian Shrinkage as an elegant, robust, and self-stabilizing regularizer that naturally bridges unregularized moment matching and uncalibrated classification heads.

\subsection{Computational and Architectural Advantages}
\label{subsec:advantages}

Compared to gradient-based calibration (REDA-SFT), MOMO-Merge provides severe practical advantages:
\begin{itemize}
  \item \textbf{Zero SGD/Adam Optimization:} No optimization algorithm, weight decay, learning rate, or number of training epochs need to be selected or tuned.
  \item \textbf{No Backpropagation:} Since it is entirely closed-form, MOMO-Merge requires only a single forward pass over the calibration set to compute statistics. This makes it computationally instantaneous (+0.0s overhead) and extremely suitable for resource-constrained edge devices or on-device model merging.
  \item \textbf{Overfitting Immunity:} Iterative gradient descent on extremely small calibration sets (e.g., $N=4$) is highly prone to overfitting, as shown by REDA-SFT's drop to 28.83\%. By relying on robust low-order statistical moments with analytical regularization, MOMO-Merge is inherently resistant to overfitting.
\end{itemize}

\section{Conclusion and Future Work}
\label{sec:conclusion}

In this work, we challenged the assumption that joint backbone and head calibration for multi-task model merging requires expensive and complex gradient-based fine-tuning. Grounded in the philosophy of Occam's razor, we introduced \textbf{MOMO-Merge}, a 100\% training-free, computationally instantaneous, and zero-hyperparameter calibration framework. MOMO-Merge combines global layer-wise activation calibration (SP-TAAC) to restore representation variance with an analytical Feature-Wise Moment Matching (FWMM) to align classification heads in closed-form. To handle extremely low-data regimes ($N \le 16$), we proposed a robust Bayesian Variance Shrinkage regularizer that stabilizes empirical statistics under ReLU sparsity. Extensive evaluations on a multi-task vision benchmark demonstrate that MOMO-Merge completely matches or exceeds gradient-based head fine-tuning while removing optimization complexity and backpropagation entirely.

In future work, we plan to extend MOMO-Merge to large language models (LLMs) and vision-language models (VLMs), exploring its application to token-level classification heads and instruction-tuning alignment.

\bibliographystyle{icml2026}
\bibliography{submission}

\newpage
\appendix
\onecolumn

\section{Mathematical Proofs and Analytical Derivations}
\label{app:proofs}

In this section, we provide the complete mathematical proof and step-by-step derivation for the closed-form absorption of the Feature-Wise Moment Matching (FWMM) transformation into a linear classification head.

Let $x \in \mathbb{R}^D$ be the intermediate representation produced by the calibrated merged backbone, and let $z^{(k)} \in \mathbb{R}^D$ be the representation produced by the original task-specific expert $k$.
The original pre-trained linear classification head for task $k$ is given by:
\begin{equation}
h_k(x) = W^{(k)} x + b^{(k)}
\end{equation}
where $W^{(k)} \in \mathbb{R}^{C \times D}$ and $b^{(k)} \in \mathbb{R}^C$.

On the task-specific calibration set of size $N$, we compute the empirical channel-wise means and standard deviations:
\begin{align}
\mu_{\text{expert}}^{(k)} &= \frac{1}{N}\sum_{i=1}^N z^{(k)}_i, & \sigma_{\text{expert}}^{(k)} &= \sqrt{\frac{1}{N}\sum_{i=1}^N (z^{(k)}_i - \mu_{\text{expert}}^{(k)})^2 + \epsilon} \\
\mu_{\text{merged}}^{(k)} &= \frac{1}{N}\sum_{i=1}^N x_i, & \sigma_{\text{merged}}^{(k)} &= \sqrt{\frac{1}{N}\sum_{i=1}^N (x_i - \mu_{\text{merged}}^{(k)})^2 + \epsilon}
\end{align}
Our goal is to apply an affine transformation to $x$ to match the first and second moments of the expert's representations $z^{(k)}$ channel-wise:
\begin{equation}
x' = (x - \mu_{\text{merged}}^{(k)}) \odot S + \mu_{\text{expert}}^{(k)}
\end{equation}
where $S = \sigma_{\text{expert}}^{(k)} / \sigma_{\text{merged}}^{(k)} \in \mathbb{R}^D$.
We can express the element-wise multiplication $\odot S$ as a matrix multiplication with a diagonal matrix $\operatorname{diag}(S) \in \mathbb{R}^{D \times D}$:
\begin{equation}
x' = \operatorname{diag}(S) (x - \mu_{\text{merged}}^{(k)}) + \mu_{\text{expert}}^{(k)}
\end{equation}
Plugging this aligned feature $x'$ directly into the pre-trained linear head $h_k(x')$:
\begin{align}
h'_k(x) &= W^{(k)} x' + b^{(k)} \\
&= W^{(k)} \left[ \operatorname{diag}(S) (x - \mu_{\text{merged}}^{(k)}) + \mu_{\text{expert}}^{(k)} \right] + b^{(k)} \\
&= W^{(k)} \operatorname{diag}(S) x - W^{(k)} \operatorname{diag}(S) \mu_{\text{merged}}^{(k)} + W^{(k)} \mu_{\text{expert}}^{(k)} + b^{(k)} \\
&= \left[ W^{(k)} \operatorname{diag}(S) \right] x + \left[ b^{(k)} + W^{(k)} \left( \mu_{\text{expert}}^{(k)} - \mu_{\text{merged}}^{(k)} \odot S \right) \right]
\end{align}
Thus, we can define the updated classification head parameters $W^{(k)\prime}$ and $b^{(k)\prime}$ as:
\begin{align}
W^{(k)\prime} &= W^{(k)} \operatorname{diag}(S) \\
b^{(k)\prime} &= b^{(k)} + W^{(k)} \left( \mu_{\text{expert}}^{(k)} - \mu_{\text{merged}}^{(k)} \odot S \right)
\end{align}
where $\operatorname{diag}(S)$ scales the columns of the weight matrix $W^{(k)}$ element-wise, and the shift is absorbed as a direct modification to the bias vector $b^{(k)}$.
At inference time, evaluating $h'_k(x) = W^{(k)\prime} x + b^{(k)\prime}$ requires exactly the same number of floating-point operations (FLOPs) and parameter storage as the original head, introducing **zero deployment or inference overhead**.

\section{Analysis of the Sparsity Trap under ReLU Activations}
\label{app:sparsity_trap}

The "sparsity trap" is a major challenge in multi-task model merging, first identified at the activation layer level by \citet{SPTAAC26}. Here, we analyze its manifestation at the classification head interface.

Under ReLU activations, a significant fraction of representation features are exactly zero (dead channels). On a very small calibration set (e.g., $N=4$ samples), many channels in the feature vector $x \in \mathbb{R}^{512}$ will be inactive across all $N$ samples. This causes the empirical standard deviation $\sigma_{\text{merged}}^{(k)}$ to be exactly zero or extremely close to the stability limit $\epsilon = 10^{-8}$.
Consequently, the unregularized empirical scaling factor $S_d = \sigma_{\text{expert}, d} / \sigma_{\text{merged}, d}$ for an inactive channel $d$ explodes to $10^{8}$, amplifying any small numerical noise in the backbone and completely corrupting the classification decision boundary.

Our proposed **MOMO-Merge with Regularized Variance Shrinkage** mathematically resolves this.
Rather than regularizing the scale $S$ after division, we regularize the standard deviations themselves by interpolating them with their global channel averages:
\begin{equation}
\sigma_{\text{merged, reg}} = (1 - \lambda) \bar{\sigma}_{\text{merged}} + \lambda \sigma_{\text{merged}}
\end{equation}
Because the global channel average $\bar{\sigma}_{\text{merged}} = \frac{1}{D}\sum_{d=1}^D \sigma_{\text{merged}, d}$ is computed across all 512 channels, it is highly robust and guaranteed to be strictly positive and large (typically $\bar{\sigma}_{\text{merged}} \approx 0.5$).
For extremely small $N=4$, $\lambda = 0.2$, which bounds the minimum standard deviation of any channel from below by:
\begin{equation}
\sigma_{\text{merged, reg}, d} \ge (1 - \lambda) \bar{\sigma}_{\text{merged}} \approx 0.8 \times 0.5 = 0.4
\end{equation}
This completely prevents division-by-zero and noise amplification, stabilizing the scales and ensuring robust decision boundaries.

\section{Computational Complexity and Wall-Clock Time Comparison}
\label{app:complexity}

In Table~\ref{tab:wall_clock}, we compare the theoretical and empirical computational complexity of MOMO-Merge versus gradient-based head calibration (REDA-SFT).

\begin{table}[h]
\caption{Wall-clock calibration time and optimization complexity on a single task with calibration size $N=32$.}
\label{tab:wall_clock}
\vskip 0.15in
\begin{center}
\begin{small}
\begin{sc}
\begin{tabular}{lccc}
\toprule
Method & Optimization Epochs & Backpropagation & Calibration Time \\
\midrule
Uncalibrated WA & 0 & No & 0.0s \\
SP-TAAC Only & 0 & No & <0.1s \\
MOMO-Merge (Ours) & 0 & No & <0.1s \\
REDA-SFT (Gradient) & 15 & Yes & ~12.5s \\
\bottomrule
\end{tabular}
\end{sc}
\end{small}
\end{center}
\vskip -0.1in
\end{table}

While REDA-SFT must run 15 epochs of iterative forward-backward passes, computing gradients and updating weights using AdamW, MOMO-Merge requires only a single forward pass over the calibration set to extract statistics, followed by an instantaneous closed-form linear algebra update. This results in a speedup of over **125$\times$** during the calibration phase, while introducing no hyperparameters to tune.

\end{document}
